\unnumberedChapter{ЗАКЛЮЧЕНИЕ}

Решены все поставленные задачи работы и достигнуты следующие основные результаты:
\begin{enumerate}
    \item \textbf{Алгоритм внутрисхемного измерения частотных характеристик}, основанный на инжекции в систему синусоидальных сигналов с последовательной сменой частоты и амплитуды, обеспечивающий оптимальный размер измеряемых сигналов-откликов системы с целью уменьшения продолжительности измерений.

    \item \textbf{Алгоритм идентификации параметров линейного динамического звена по его частотной характеристике}, основанный на архитектуре энкодера-декодера, обеспечивающей однократную генерацию узких бинарных масок, соответствующих координатам частот вещественных полюсов и нулей передаточной функции звена.
\end{enumerate}

В случае высокоскоростного канала связи, когда третье слагаемое выражения \eqref{eqn:meas_duration} уже не преобладает в полной сумме, на передний план выходит задача сокращения длительности измерительного окна, т.е. уменьшение параметра $L$. Обозначенная задача отнесена к перспективе дальнейших исследований.

Важно отметить, что архитектура модели сегментации разработана под широкий класс задач автоматизированного контроля и тестирования встраиваемых систем управления. Рассмотренная в данной работе задача детекции полюсов/нулей относится к простым примерам применения модели, под которые возможно синтезировать обучающий набор данных. Модель возможно обучить детекции разного рода отклонений при наличии соответствующих наборов данных.